\documentclass[11pt,leqno]{article}
\usepackage[margin=1in]{geometry}
\usepackage{amssymb}
\usepackage{amsmath}
\usepackage{amsthm}
\usepackage{booktabs}
\usepackage{verbatim}

\renewcommand{\labelenumi}{(\arabic{enumi})}
\emergencystretch=3em

\newtheorem{theorem}[equation]{Theorem}
\newtheorem{corollary}[equation]{Corollary}
\newtheorem{lemma}[equation]{Lemma}
\newtheorem{proposition}[equation]{Proposition}
\theoremstyle{definition}
\newtheorem{definition}[equation]{Definition}
\newtheorem{observation}[equation]{Observation}
\newtheorem{heuristic}[equation]{Working hypothesis}
\newtheorem{remark}[equation]{Remark}
\newtheorem{example}[equation]{Example}
\numberwithin{equation}{section}

\newcommand{\K}{\mathrm{K}}
\newcommand{\nil}{\mathcal N}
\newcommand{\Ob}{\mathcal O}
\newcommand{\fg}{\mathfrak g}
\newcommand{\fk}{\mathfrak k}
\newcommand{\fs}{\mathfrak s}
\newcommand{\fq}{\mathfrak q}
\newcommand{\fu}{\mathfrak u}
\newcommand{\fl}{\mathfrak l}
\newcommand{\ft}{\mathfrak t}
\newcommand{\Z}{\mathbb Z}
\newcommand{\Q}{\mathbb Q}
\newcommand{\R}{\mathbb R}
\newcommand{\C}{\mathbb C}
\newcommand{\ch}[1]{#1^\vee}
\newcommand{\AV}{\mathrm{AV}}
\newcommand{\AC}{\mathrm{AC}}
\newcommand{\dom}{\mathrm{dom}}
\newcommand{\at}[1]{\texttt{#1}}
\newcommand{\mi}{\medskip\noindent}

\begin{document}
\title{The $K$-nilpotent algorithm with smaller matrices}
\author{Draft prepared with Claude for J.~Adams}
\date{November 2026}
\maketitle

These notes propose a modification of the $K$-nilpotent algorithm of
\cite{AV}, as implemented in \at{K\_Nilpotent.at} and described in
\at{knilpotent.tex} (referred to below as \cite{KN}).  The notation is that of
\cite{KN}.  Everything here was tested in \at{atlas}; the test code is in
Appendix~\ref{s:code}.  Statements are labelled according to their status:
\emph{Theorem}, \emph{Proposition} and \emph{Lemma} are proved (a proof
or a sketch is given); an \emph{Observation} is an experimental fact;
a \emph{Working hypothesis} is a rule that the algorithm tries first, and
which it does not need to be true.

\section{Summary}
\label{s:summary}

The current implementation fixes a bound $N$, computes $\Phi(\Ob,\sigma)$ for
\emph{every} real nilpotent orbit $\Ob$ and every $L\cap K$-type $\sigma$ of
norm at most $2N$, keeps those all of whose terms have $K$-norm at most $N$,
and then does lattice computations (intersections, quotients) with the
matrices $T$, $Y$, $Q$, $P_X$ of \cite[Section~12]{KN}.  Measurements
(Section~\ref{s:cost}) show that the rows of these matrices are not the
problem; the columns are.  There are four to twelve times as many columns as
rows, and eighty to ninety percent of the $\Phi(\Ob,\sigma)$ that are
computed are thrown away.

The proposal consists of four changes.

\begin{enumerate}
\item \textbf{Read the answer off any solution} (Section~\ref{s:mult}).
If
$[\pi]=\sum a_{\Ob,\sigma}\,\Phi(\Ob,\sigma)$ is \emph{any} expression of
$[\pi]$ in terms of orbits of dimension at most $d$, with rational
coefficients, then the multiplicity of $\overline\Ob$ in the associated cycle
of $\pi$ is $\sum_\sigma a_{\Ob,\sigma}\dim\sigma$ for every orbit of
dimension $d$ (Theorem~\ref{t:mult}).  So there is no need for the matrices
$Y$, $Q$ and $P_X$, no need for integer (Smith form) linear algebra, and no
need to know that all relevant $\sigma$ have been found: one linear system
$Tv=w$ over $\Q$ is solved, and the answer certifies itself.  As a bonus the
algorithm returns the multiplicities, not only the associated variety.

\item \textbf{Use only the orbits of the right dimension}
(Section~\ref{s:top}).  The Gelfand--Kirillov dimension of $\pi$ is available
from the $W$-cell of $\pi$ (\at{GK\_dim}).  In most cases $[\pi]$ is a
combination of the $\Phi(\Ob,\sigma)$ with $\dim\Ob=2\,\mathrm{GKdim}(\pi)$
alone.  When it is not, smaller orbits are added, largest first.

\item \textbf{The cone argument} (Section~\ref{s:cone}).  Order the $K$-types
by \emph{dominance} of the weight $\gamma(\mu)=(1+\theta_x)\lambda$ rather
than by the single number $\|\mu\|=\langle\gamma(\mu),\ch\rho\rangle$.  The
$K$-type formula, branching to $K$, standardization, and $\Phi$ are all
triangular for this finer order.  Only the $K$-types lying between the
lowest and the highest terms of $[\pi]$ are used (the intersection of an
upward cone and a downward cone), and a $\sigma$ is discarded, before
$\Phi(\Ob,\sigma)$ is computed, if the leading term of $\Phi(\Ob,\sigma)$ is
outside this region.

\item \textbf{An additive bound on $\sigma$} (Section~\ref{s:bound}).  The
$\sigma$ that can matter have norm at most $N+c_\Ob$ for an explicit constant
$c_\Ob$, in place of the bound $2N$ of the current code.  This is the form of
the bound in \cite[Proposition~8.9]{AV}.  The bound $2N$ is wasteful when $N$
is large, and \emph{too small} when $N$ is small: see Example~\ref{e:wrong},
where the current code returns a wrong answer.
\end{enumerate}

Because of (1), the rules in (2)--(4) do not have to be theorems.  They decide
which columns are tried first.  If the linear system has a solution, the
answer is correct; if not, the rules are relaxed and the computation repeated
with more columns.

The effect is summarized in Table~\ref{tab:blocks} of Section~\ref{s:exp}.
For the block of the trivial representation of $SU(2,2)$ the time goes from
81 seconds to 3.2 seconds, for $SO(4,3)$ from 136 to 34, and for
$SL(5,\R)$ from 68 to 11.  For $Sp(6,\R)$ the current code did not finish
in 50 minutes; the new one takes 51 seconds.  For a single representation
the gain is larger: a large discrete series of $F_4$, $Sp(8,\R)$ or $SO(5,4)$
leads to a linear system of size at most $2\times 2$.

\section{The current algorithm, and where the time goes}
\label{s:cost}

We recall the setting briefly.  Let $M=\K^K(\nil_\theta)$.  It has the basis
$\{[\mu]\}$ indexed by $K$-types, $[\mu]$ being the class of the standard
tempered final limit module with lowest $K$-type $\mu$.  For a real nilpotent
orbit $\Ob=K\cdot E$, with $\theta$-stable parabolic $Q=LU$ defined by the
semisimple element $H$ of a normal triple, there is the map
\[
\Phi(\Ob,\cdot)\colon R(L\cap K)\longrightarrow M ,
\]
computed by the formula of \cite[Section~6]{KN}.  The \emph{dimension} of a real orbit
always means the dimension of the complex orbit $G\cdot\Ob$, as in the
software (it is twice the dimension of $\Ob$ as a $K$-orbit).  The class of an
irreducible $\pi=J(p)$ is $w=[\pi]$, given in the basis $\{[\mu]\}$ by
\at{character\_formula(p)*0}.

The current algorithm sets $N=\|w\|$, lets $W(N)$ be the span of the $[\mu]$
with $\|\mu\|\le N$, and computes for every real orbit all $\sigma$ with
$\|\sigma\|\le 2N$ and $\Phi(\Ob,\sigma)\in W(N)$.

Table~\ref{tab:old} records, for the block of the trivial representation, the
number $m$ of $K$-types of norm at most $N$ (the number of rows), the number
of $\Phi(\Ob,\sigma)$ that are computed, and the number that are kept (the
number of columns).

\begin{table}[ht]
\centering
\begin{tabular}{lrrrrrr}
\toprule
$G$ & block & $N$ & rows $m$ & $\Phi$ computed & $\Phi$ kept & time (s)\\
\midrule
$SL(2,\R)$ & 3 & 1 & 5 & 17 & 7 & 0.01\\
$SU(2,1)$ & 6 & 4 & 16 & 226 & 48 & 0.2\\
$SL(3,\R)$ & 6 & 4 & 4 & 24 & 10 & 0.03\\
$Sp(4,\R)$ & 12 & 7 & 25 & 522 & 97 & 1.0\\
$G_2$ split & 12 & 16 & 16 & 558 & 105 & 1.8\\
$SL(4,\R)$ & 15 & 10 & 21 & 824 & 137 & 3.2\\
$SU(2,2)$ & 21 & 10 & 95 & 7395 & 698 & 81\\
$SO(4,3)$ & 53 & 22 & 55 & 6564 & 612 & 136\\
$SL(5,\R)$ & 55 & 20 & 21 & 1390 & 240 & 68\\
$Sp(6,\R)$ & 53 & 22 & 146 & 29083 &  & $>3000$\\
$SU(3,2)$ & 55 & 20 & 413 & 221514 &  & $>3000$\\
\bottomrule
\end{tabular}
\caption{The current algorithm on the block of the trivial representation.
The time is for filling the data and computing the associated variety of
every irreducible in the block. For the last two groups the computation was
stopped after 50 minutes.}
\label{tab:old}
\end{table}

Three things stand out.

\begin{itemize}
\item The number of rows is modest.  Even for split $F_4$ there are only 606
$K$-types of norm at most $N=\langle 2\rho,\ch\rho\rangle=110$, and for
$Sp(8,\R)$ there are 973.
\item Most of the work is computing $\Phi(\Ob,\sigma)$ for $\sigma$ that are
then discarded because some term has norm bigger than $N$.
\item The columns that are kept are highly redundant.  For the principal orbit
of split $G_2$ there are 33 columns, and the matrix $Q$ has one row; for the
principal orbit of $SO(4,3)$ there are 165 columns and $Q$ has two rows.
\end{itemize}

\section{Multiplicities from any solution}
\label{s:mult}

For $d\ge 0$ let $Z_d\subset\nil_\theta$ be the union of the $K$-orbits of
dimension at most $d$.  It is closed.  By \cite[Theorem~4.5]{AV}
the map $\K^K(Z_d)\to\K^K(\nil_\theta)$ is injective, with free
cokernel; write $F_d\subset M$ for its image.  By
\cite[Corollary~12.2 and Algorithm~12.4]{AV}, $F_d$ is the span of the
$\Phi(\Ob,\sigma)$ with $\dim\Ob\le d$.

\begin{lemma}
\label{l:rank}
Let $\Ob=K\cdot E$ be an orbit of dimension $d$.  There is a homomorphism
\[
m_\Ob\colon F_d\longrightarrow \Z
\]
with the following properties.
\begin{enumerate}
\item $m_\Ob(\Phi(\Ob,\sigma))=\dim\sigma$ for every virtual representation
$\sigma$ of $L\cap K$.
\item $m_\Ob(\Phi(\Ob',\sigma'))=0$ for every orbit $\Ob'\ne\Ob$ of dimension
at most $d$.
\item If $X$ is a $(\fg,K)$-module of finite length with $\dim\AV(X)\le d$,
then $[X]\in F_d$, and $m_\Ob([X])$ is the multiplicity of $\overline\Ob$ in
the associated cycle of $X$.  It is nonnegative, and it is positive if and
only if $\overline\Ob$ is an irreducible component of $\AV(X)$.
\end{enumerate}
\end{lemma}

\begin{proof}
The orbit $\Ob$ is open in $Z_d$.  Restriction to an open set is exact, so it
defines $\K^K(Z_d)\to\K^K(\Ob)\simeq R(K^E)$; let $m_\Ob$ be the composition
of this map with the dimension $R(K^E)\to\Z$.  Statement~(1) is
\cite[(12.4d)]{AV}: the map $\mu\colon K\times_{Q\cap K}\fs[\ge 2]\to
\overline\Ob$ is an isomorphism over $\Ob$, and restriction to $\Ob$ commutes
with $\mu_*$.  For (2), $\Phi(\Ob',\sigma')$ is supported on $\overline{\Ob'}$,
which does not meet $\Ob$.  Statement~(3) is the definition of the
associated cycle \cite[Definition~5.2]{AV}: $m_\Ob([X])$ is the rank of the
coherent sheaf $\mathrm{gr}X$ at a point of $\Ob$.
\end{proof}

\begin{theorem}
\label{t:mult}
Let $X$ be a $(\fg,K)$-module of finite length.  Suppose that for some $d$
\begin{equation}
\label{e:expr}
[X]=\sum_{\dim\Ob'\le d}\ \sum_{\sigma}a_{\Ob',\sigma}\,\Phi(\Ob',\sigma)
\qquad(a_{\Ob',\sigma}\in\Q),
\end{equation}
a finite sum.  For each orbit $\Ob$ of dimension $d$ set
\[
m_\Ob=\sum_\sigma a_{\Ob,\sigma}\dim\sigma .
\]
\begin{enumerate}
\item $\dim\AV(X)\le d$, and $m_\Ob$ is the multiplicity of $\overline\Ob$ in
the associated cycle of $X$.  In particular $m_\Ob$ is a nonnegative integer
which does not depend on the expression \eqref{e:expr}.
\item If $m_\Ob\ne0$ for some $\Ob$ of dimension $d$, then $\dim\AV(X)=d$ and
the components of $\AV(X)$ of dimension $d$ are the $\overline\Ob$ with
$m_\Ob\ne 0$.  If $X$ is irreducible these are all the components of
$\AV(X)$.
\item If $m_\Ob=0$ for all $\Ob$ of dimension $d$, then $\dim\AV(X)<d$.
\end{enumerate}
\end{theorem}

\begin{proof}
By \eqref{e:expr}, $[X]\in F_d\otimes\Q$, hence $[X]\in F_d$ because $M/F_d$
is free.  Suppose $\AV(X)$ had a component $\overline{\Ob''}$ of dimension
$d''>d$, with $d''$ maximal.  Then $m_{\Ob''}$ is defined on $F_{d''}$ and is
positive on $[X]$ by Lemma~\ref{l:rank}(3), but it vanishes on
$F_{d''-1}\supset F_d$ by Lemma~\ref{l:rank}(2).  So $\dim\AV(X)\le d$.  Now
apply $m_\Ob\otimes\Q$ to \eqref{e:expr} and use Lemma~\ref{l:rank}.  This
gives (1), and (2) and (3) follow from (1) and Lemma~\ref{l:rank}(3).  The last
statement of (2) is \cite[Theorem~8.4]{Vogan91}: the components of the
associated variety of an irreducible module all have the same dimension.
\end{proof}

\begin{remark}
\label{r:consequences}
Here is what Theorem~\ref{t:mult} does for the algorithm.

\smallskip\noindent
(a) The matrices $Y$, $Q$, $P_X$ of \cite[Section~12]{KN}, and the
computation of \at{intersection} and \at{mod\_image\_projector}, are not
needed.  One solves $Tv=w$, where the columns of $T$ are whatever
$\Phi(\Ob',\sigma)$ have been computed.

\smallskip\noindent
(b) Rational coefficients suffice, so $Tv=w$ can be solved by Gaussian
elimination (\at{solve} with a \at{ratvec} right hand side).

\smallskip\noindent
(c) The conjecture in \cite[Section~11]{KN} that the spaces $Z(\Ob_i,N)$ are
linearly independent is not needed, and neither is the completeness of
$V(\Ob,N)$ and $Y(\Ob,N)$.  \emph{Any} set of columns for which $Tv=w$ has a
solution gives the right answer.  Which columns to compute becomes a question
of efficiency only.

\smallskip\noindent
(d) The test ``the block of $v$ belonging to $\Ob_i$ is nonzero'' in
\at{av\_from\_av\_ann} is replaced by ``$m_{\Ob_i}\ne0$''.  These are different.
For the irreducible $J(x=10,\lambda=[3,2],\nu=[2,1])$ of $Sp(4,\R)$ the
smallest solution we found uses columns from all three real forms of the
orbit of dimension $6$, and the multiplicities are $(0,0,1)$.

\smallskip\noindent
(e) The answer checks itself: the $m_\Ob$ must be nonnegative integers, not
all zero.  If they all vanish, the true associated variety is smaller, and
the columns used for the smaller orbits were insufficient.
\end{remark}

\begin{example}
\label{e:wrong}
Let $G=Sp(4,\R)$ and let $\pi=J(p)$ be the limit of holomorphic discrete
series with $x=3$ and $\lambda=(1,0)$, at the singular infinitesimal character
$(1,0)$.  Here $w=[\pi]$ is a single term, of norm $N=3$, and
\at{GK\_dim(p)}$=3$, so $\AV(\pi)$ has dimension $6$.  With $N=3$ the current
code gives
\begin{verbatim}
atlas> set p=parameter(KGB(G,3),[1,0],[0,0])
atlas> set d=fill(initialize_KNilpotent(G,3))
atlas> av(d,p)
Value: [ 7 ]
\end{verbatim}
which is a principal orbit, of dimension $8$.  With $N\ge 5$ it gives
$[4]$, which is correct.  What happens with $N=3$ is that $w$ is not in the
span of the columns that were computed for the orbits of dimension $6$,
because the needed $\sigma$ have norm $9$, which is bigger than $2N=6$
(see Section~\ref{s:bound}); but it is in the span once the columns of the
principal orbits are added.  The multiplicities
of that solution are $(0,0)$, so Theorem~\ref{t:mult}(3) detects the problem.
So the statement in \cite[Section~10]{KN} that one may take $N=\|P\|_v$ is
not safe, and with the current code there is no way to tell from the output
whether $N$ was big enough.  For $SL(2,\R)$ at infinitesimal character $0$
one has $N=0$, and the current code returns nothing for the two limits of
discrete series and for the spherical principal series.
\end{example}

\section{Which orbits}
\label{s:top}

Let $\pi$ be irreducible and $d=2\,\mathrm{GKdim}(\pi)$.  The function
\at{GK\_dim} (in \at{associated\_variety\_}\allowbreak\at{annihilator.at}) computes
$\mathrm{GKdim}(\pi)$ from the special representation of the $W$-cell of
$\pi$; given the Kazhdan--Lusztig data that \at{character\_formula} needs
anyway, this costs nothing.  So the search of \cite[Section~10.1]{KN} over
complex orbits in increasing order can be skipped, and no orbit of dimension
greater than $d$ is ever needed.

\begin{observation}
\label{o:top}
In most cases $[\pi]$ is a rational combination of the $\Phi(\Ob,\sigma)$ with
$\dim\Ob=d$ exactly; no smaller orbits are needed.
\end{observation}

For the block of the trivial representation this holds for every irreducible
for $SL(2,\R)$, $SU(2,1)$, $SL(3,\R)$, $Sp(4,\R)$, split $G_2$, $SL(4,\R)$,
$SU(2,2)$ and $Sp(6,\R)$, for 47 of the 53 irreducibles for $SO(4,3)$, and
for 43 of the 55 for $SL(5,\R)$: in all, for 218 of 236
(Table~\ref{tab:blocks}).

It is not a theorem.  The image of $\Phi(\Ob,\cdot)$ is the image of
\[
\mu_*\colon \K^K(K\times_{Q\cap K}\fs[\ge 2])\longrightarrow
\K^K(\overline\Ob).
\]
This maps onto $\K^K(\overline\Ob)/\K^K(\partial\Ob)$, but it need not contain
$\K^K(\partial\Ob)$: over a boundary orbit $\Ob'=K\cdot E'$ the fibre of $\mu$
is a projective variety on which $K^{E'}$ acts, and $\mu_*$ produces only the
Euler characteristics of equivariant sheaves on that fibre.

\begin{example}
\label{e:so43}
Let $G=SO(4,3)$.  The complex orbit with partition $[3,3,1]$ has two real
forms, numbered 5 and 6 by \at{real\_nilpotent\_orbits(G)}.  There are six
irreducibles in the block of the trivial representation whose associated
variety is $\overline{\Ob_5}\cup\overline{\Ob_6}$ (with multiplicities
$(1,1)$, $(4,4)$ and $(5,5)$).  For these, $[\pi]$ is not in the span of the
$\Phi(\Ob_5,\sigma)$ and $\Phi(\Ob_6,\sigma)$.  It is in the span after the
columns of orbit number 2 are added; this is a real form of the orbit
$[3,1^4]$, of dimension 10.  For every irreducible in the block whose
associated variety is $\overline{\Ob_5}$ alone, or $\overline{\Ob_6}$ alone,
the top orbits suffice.  These six are the only irreducibles in the block
with a reducible associated variety, and the only ones for which the top
orbits do not suffice.

A reducible associated variety is not the only cause.  In the block of the
trivial representation of $SL(5,\R)$ there are twelve irreducibles, each with
an irreducible associated variety, for which the orbits of the next smaller
dimension are needed.
\end{example}

So the rule is: try the orbits of dimension $d$; if $Tv=w$ has no solution,
add the orbits of the next smaller dimension, and so on.  By
Theorem~\ref{t:mult} nothing is lost by trying.  It would be better to add
only the orbits in the closure of at least two of the orbits found so far,
but this requires the closure order on real nilpotent orbits.

\section{The cone argument}
\label{s:cone}

\subsection{The dominance order on $K$-types}

\begin{definition}
For a $K$-type $\mu=(x,\lambda)$ let
\[
\gamma(\mu)=\dom\big((1+\theta_x)\lambda\big)\in X^*\otimes\Q,
\]
the dominant element of the $W$-orbit of $(1+\theta_x)\lambda$.  For dominant
$\gamma,\gamma'$ write $\gamma\preceq\gamma'$ if $\gamma'-\gamma$ is a
nonnegative rational combination of simple roots.
\end{definition}

The $K$-norm of \cite[Section~8]{KN} is
$\|\mu\|=\langle\gamma(\mu),\ch\rho\rangle$, and more generally
$\|\mu\|_v=\langle\gamma(\mu),v\rangle$.  Since $\gamma\preceq\gamma'$ if and
only if $\langle\gamma,v\rangle\le\langle\gamma',v\rangle$ for all dominant
$v$ (and equality for central $v$), the dominance order uses all the norms
$\|\cdot\|_v$ at once.  For a dominant weight $\Lambda$,
\begin{align*}
\{\gamma\ \text{dominant}\mid\gamma\preceq\Lambda\}
&=\{\text{dominant weights}\}\cap\big(\Lambda-\mathrm{cone}(\Delta^+)\big)\\
&=\{\text{dominant weights in the convex hull of }W\Lambda\}.
\end{align*}
This is a cone condition.  It has two advantages over $\|\mu\|\le N$.  It is
sharper: the simplex $\{\gamma\mid\langle\gamma,\ch\rho\rangle\le N\}$ is
replaced by a polytope with vertex $\Lambda$.  And it is finite even when $G$
has a compact central torus, since the central character is fixed; no
auxiliary vector $v$ is needed.

For $w=\sum_\mu c_\mu[\mu]\in M$ let $\Gamma_{\max}(w)$ and
$\Gamma_{\min}(w)$ be the sets of maximal and minimal elements of
$\{\gamma(\mu)\mid c_\mu\ne 0\}$, and define the \emph{cone} and the
\emph{interval} of $w$:
\begin{align*}
C(w)&=\{\mu\mid\gamma(\mu)\preceq\gamma_+\text{ for some
}\gamma_+\in\Gamma_{\max}(w)\},\\
I(w)&=\{\mu\in C(w)\mid\gamma_-\preceq\gamma(\mu)\text{ for some
}\gamma_-\in\Gamma_{\min}(w)\}.
\end{align*}
The lowest $K$-types of $\pi$ give elements of $\Gamma_{\min}([\pi])$.

Table~\ref{tab:cone} compares the number of $K$-types of norm at most
$N=\langle 2\rho,\ch\rho\rangle$ with the number in the cone below $2\rho$
(which is $\Gamma_{\max}([\pi])$ for most $\pi$ of infinitesimal character
$\rho$).  The cone by itself removes between five and fifty percent of the
rows.  The
interval removes many more when the lowest $K$-type of $\pi$ is large: for a
discrete series of $Sp(4,\R)$ the interval contains 4 of the 25 $K$-types, for
split $G_2$ it contains 3 of 16, and for a fundamental series of $SL(4,\R)$ it
contains 2 of 21.

\begin{table}[ht]
\centering
\begin{tabular}{lrrr}
\toprule
$G$ & $N$ & $\|\mu\|\le N$ & $\gamma(\mu)\preceq 2\rho$\\
\midrule
$SU(2,2)$ & 10 & 95 & 63\\
$SO(4,3)$ & 22 & 55 & 52\\
$Sp(6,\R)$ & 22 & 146 & 132\\
$SU(3,2)$ & 20 & 413 & 209\\
$SO(5,4)$ & 50 & 328 & 264\\
$Sp(8,\R)$ & 50 & 973 & 731\\
$F_4$ split & 110 & 606 & 544\\
\bottomrule
\end{tabular}
\caption{Number of $K$-types of norm at most $N$, and in the cone below
$2\rho$.}
\label{tab:cone}
\end{table}

\subsection{Triangularity}

The reason to expect the cone and the interval to be respected is that every
operation that enters the computation of $[\pi]$ and of $\Phi(\Ob,\sigma)$ is
triangular with respect to $\preceq$.  The current algorithm uses only the
consequence of this for the single functional
$\langle\,\cdot\,,\ch\rho\rangle$.

\begin{observation}[$K$-type formula and branching]
\label{o:branch}
Every term $[\mu']$ of the $K$-type formula of $\mu$ satisfies
$\gamma(\mu)\preceq\gamma(\mu')$.  Every $K$-type $\mu'$ of the standard module
with lowest $K$-type $\mu$ satisfies $\gamma(\mu)\preceq\gamma(\mu')$.
\end{observation}

This was checked on all $K$-types up to a bound for the groups $SL(2,\R)$,
$SU(2,1)$, $SL(3,\R)$, $Sp(4,\R)$, $G_2$, $SL(4,\R)$, $SU(2,2)$, $SO(4,3)$,
$Sp(6,\R)$: about 5500 terms of $K$-type formulas and 5000 branching
multiplicities, with no exception.  It refines the fact that
\at{K\_type\_formula} and \at{branch} are triangular for the height.  We expect
that it follows from Vogan's theory of lowest $K$-types, since
$\gamma(\mu)=2\lambda_a(\mu)$, but we have not written a proof.

\begin{lemma}[standardization]
\label{l:std}
Let $x$ be a KGB element and $\lambda\in\rho+X^*$, not necessarily standard
or final, and let $[I(x,\lambda)]=\sum_j c_j[\mu_j]$ be the expression of the
continued standard module in the basis of $M$.  Then
$\gamma(\mu_j)\preceq\dom((1+\theta_x)\lambda)$ for all $j$.
\end{lemma}

\begin{proof}[Sketch of proof]
The set of $\gamma$ with $\dom(\gamma)\preceq\Lambda$ is the convex hull of
$W\Lambda$, so it is convex and $W$-stable.  The parameter is made standard
and final by a sequence of steps of four kinds.  Reflection in a compact
imaginary simple root, and the cross action of a complex root, replace
$(1+\theta_x)\lambda$ by a $W$-conjugate.  The Hecht--Schmid identity for a
noncompact imaginary simple root $\alpha$ replaces $(x,\lambda)$ by terms
at $s_\alpha\times x$, where $(1+\theta)\lambda$ is replaced by its reflection,
and at the Cayley transform of $x$, where it is replaced by
$(1+\theta_x)\lambda-\langle\lambda,\ch\alpha\rangle\alpha$, the midpoint of
$(1+\theta_x)\lambda$ and its reflection in $\alpha$.  Finally the inverse
Cayley transform through a real root satisfying the parity condition (the
step that makes a parameter final) does not change $(1+\theta)\lambda$.
\end{proof}

Lemma~\ref{l:std} was also tested: for every $\sigma$ of norm at most $2N$
for $Sp(4,\R)$, $SL(4,\R)$, $G_2$ and $SU(2,1)$ (about 2000 in all), whenever
the terms of $\Phi(\Ob,\sigma)$ before standardization all lie below
$2\rho$, so do the terms after standardization.

The third statement is about $\Phi$, and it is the one that says that each
orbit has its own cone.  Fix a real orbit $\Ob=K\cdot E$ with normal triple
$(H,E,F)$, $H\in\ft\subset\fk$, and $\fq=\fl+\fu$.  Choose a positive
system for $\ft$ in $\fk$ containing the roots of $\fu\cap\fk$, so that $H$ is
dominant for $K$.  For a $K$-type with highest weight $\gamma$ write
\[
h_\Ob(\gamma)=\langle\gamma,H\rangle .
\]
The element $H$ is central in $\fl\cap\fk$, so it acts on an irreducible
representation $\sigma$ of $L\cap K$ by a scalar, written $\sigma(H)$.

\begin{proposition}
\label{p:cone}
Let $\sigma$ be an irreducible representation of $L\cap K$ with highest
weight $\tau$.
\begin{enumerate}
\item As a virtual representation of $K$,
\[
\Phi(\Ob,\sigma)|_K=\sum_{k\ge 0}\ \sum_i(-1)^i\,
\mathcal R^i_{\fq\cap\fk}\big(\sigma\otimes S^k(\fs[\ge 2])\big),
\]
where $\mathcal R^i_{\fq\cap\fk}$ is cohomological induction from $L\cap K$ to
$K$, normalized so that the irreducible representation of $L\cap K$ of
highest weight $\nu$ goes to $\pm$ the irreducible representation of $K$ of
highest weight $w(\nu+\rho_K)-\rho_K$, or to zero.
\item Every $K$-type of $\Phi(\Ob,\sigma)$, of highest weight $\gamma$,
satisfies $h_\Ob(\gamma)\ge\sigma(H)$.
\item Equality holds in (2) for at most one $K$-type.  It holds if and only
if $\tau$ is dominant for $K$ and $\gamma=\tau$, and then the multiplicity is
one.
\end{enumerate}
\end{proposition}

\begin{proof}
By the definition of $\Phi$ in \cite[Section~6]{KN}, in the form in which it is implemented (see Remark~\ref{r:sign}),
\[
\Phi(\Ob,\sigma)=\sum_j(-1)^j\big[\mathcal R\big(Q,\sigma\otimes
\textstyle\bigwedge^j\fs[1]\otimes\bigwedge^{\rm top}(\fu\cap\fs)^*\big)\big].
\]
By Zuckerman's Blattner formula \cite[Proposition~11.7]{AV}, with the
normalization of cohomological induction used in \at{atlas}, the restriction
to $K$ of $\sum_i(-1)^i\mathcal R^i_\fq(Z)$ is
$\sum_i(-1)^i\mathcal R^i_{\fq\cap\fk}(Z\otimes\bigwedge^{\rm
top}(\fu\cap\fs)\otimes S(\fu\cap\fs))$.  Now
$\fu\cap\fs=\fs[1]\oplus\fs[\ge 2]$ and
$\sum_j(-1)^j\bigwedge^j\fs[1]\otimes S(\fs[1])=1$, which gives (1).

Let $\nu$ be the highest weight of a constituent of
$\sigma\otimes S^k(\fs[\ge 2])$.  Then $\langle\nu,H\rangle\ge\sigma(H)+2k$.
If this constituent contributes a $K$-type, its highest weight is
$\gamma=w(\nu+\rho_K)-\rho_K$ with $w(\nu+\rho_K)$ dominant for $K$.  Since
$H$ is dominant, $\langle w(\nu+\rho_K),H\rangle\ge\langle\nu+\rho_K,H\rangle$,
so $h_\Ob(\gamma)\ge\langle\nu,H\rangle\ge\sigma(H)+2k$.  This is (2).

If equality holds then $k=0$, $\nu=\tau$, and
$\beta=w(\tau+\rho_K)-(\tau+\rho_K)$ is a nonnegative combination of positive
roots of $\fk$ which is orthogonal to $H$, hence a combination of positive
roots of $\fl\cap\fk$.  As $\tau+\rho_K$ is dominant for $L\cap K$ we have
$\langle\tau+\rho_K,\beta\rangle\ge 0$, and
$|\tau+\rho_K+\beta|=|\tau+\rho_K|$ forces $\beta=0$.  So $\tau+\rho_K$ is
dominant and regular for $K$, that is, $\tau$ is dominant for $K$, and
$\gamma=\tau$.  The converse is clear, and this is (3).
\end{proof}

Proposition~\ref{p:cone} was tested with \at{branch} on every column in the
data for $SL(2,\R)$, $SL(3,\R)$, $SU(2,1)$, $Sp(4,\R)$ and $G_2$: among about
10000 $K$-types of the various $\Phi(\Ob,\sigma)$ none has
$h_\Ob(\gamma)<\sigma(H)$, and no $\sigma$ has two $K$-types with
$h_\Ob(\gamma)=\sigma(H)$.  This also confirms the sign conventions.

\begin{remark}
\label{r:sign}
The formula in \cite[Section~6]{KN} has $-2\rho(A)$, which corresponds to
$\bigwedge^j\fs[1]^*$; the code (\at{Phi} in \at{K\_Nilpotent.at}) adds
$2\rho(A)$, which corresponds to $\bigwedge^j\fs[1]$.
Proposition~\ref{p:cone} is stated, and was tested, for the formula as it is
implemented.  With $\bigwedge^j\fs[1]^*$ the cancellation against
$S(\fs[1])$ in the proof would not take place, and the lowest $K$-type of
$\Phi(\Ob,\sigma)$ would be shifted by $-2\rho(\fs[1])$ when $\Ob$ is not
even.  This is the point about signs raised in \cite[Section~6]{KN}.
\end{remark}

So the $K$-types of $\Phi(\Ob,\sigma)$ lie in the half space
$h_\Ob\ge\sigma(H)$, more precisely in the cone with vertex $\tau$ spanned by
the weights of $\fs[\ge 2]$, folded into the dominant chamber; and when $\tau$
is dominant for $K$ the vertex is itself a $K$-type of $\Phi(\Ob,\sigma)$,
the unique one at the lowest level of $h_\Ob$.

\subsection{The restricted system}

\begin{heuristic}
\label{h:interval}
$[\pi]$ is a rational combination of those $\Phi(\Ob,\sigma)$, $\dim\Ob=d$,
all of whose terms lie in the interval $I([\pi])$.
\end{heuristic}

The reasoning is the usual one for triangular systems.  By
Observation~\ref{o:branch} the $K$-types of $\pi$ lie above
$\Gamma_{\min}([\pi])$, and by Proposition~\ref{p:cone} the $K$-types of a
column lie above its vertex.  So the lowest terms of $[\pi]$ should be
accounted for by columns that start there; after subtracting those columns
the same holds for the remainder; and the process should not need to go above
the top of $[\pi]$.  This is not a proof, because columns can
cancel against each other outside the interval, and
Working hypothesis~\ref{h:interval} is not always true.  It does not have to be:
if the restricted system has a solution, Theorem~\ref{t:mult} applies.

The evidence is as follows.  Of the 236 irreducibles in the blocks of
Table~\ref{tab:blocks}, the top orbits suffice for 218.  For 212 of these the
system restricted to the interval $I([\pi])$ has a solution (stages 0--2 of
Section~\ref{s:alg}); for the other 6 the interval is too small, and the
system restricted to the cone $C([\pi])$ has a solution.  For the 18 that need
smaller orbits, the columns of all the orbits used can again be restricted to
the cone $C([\pi])$.

When the infinitesimal character is not $\rho$ the cone can be too small.

\begin{example}
\label{e:sl2}
Let $G=SL(2,\R)$, with KGB elements $x=0,1$ for the compact Cartan and $x=2$
for the split one.  The spherical principal series at $\nu=0$ is
$w=[(2,1)]$, a single term with $\gamma=0$.  It is
\[
w=\Phi(\Ob_1,\sigma_0)+\Phi(\Ob_2,\sigma_2),\qquad
\Phi(\Ob_1,\sigma_0)=[(2,1)]-[(1,1)],\quad\Phi(\Ob_2,\sigma_2)=[(1,1)].
\]
The term $[(1,1)]$, a discrete series, has $\gamma=2\ne0$.  So the columns
must be allowed to reach above $\Gamma_{\max}(w)$.  The same happens for any
irreducible principal series, for example the spherical principal series of
$Sp(4,\R)$ at $\nu=2\rho$.
\end{example}

For this reason the cone is enlarged when necessary: $C_k(w)$ is defined like
$C(w)$ with $\gamma_+$ replaced by $\gamma_++2k\rho$.  In all the cases tried
$k\le 1$ was enough.  At infinitesimal character $\rho$ the terms of $[\pi]$
on the fundamental Cartan already have $\gamma=2\rho$, the largest value that
can occur for the standard modules of the block, and $k=0$ has always been
enough.

\subsection{Choosing $\sigma$ before computing $\Phi(\Ob,\sigma)$}
\label{s:bound}

Computing $\Phi(\Ob,\sigma)$ is expensive: it requires the $K$-type formula
of $\sigma$ on $L$, a sum over subsets of the weights of $\fs[1]$, and a
standardization on $G$ for every term.  So one wants to know \emph{before}
computing it whether its terms can lie in the region $R$ (an interval or a
cone) under consideration.

\begin{definition}
Let $\sigma$ be an $L\cap K$-type, with standard final limit parameter
$(x_L,\lambda_L)$ for $L$.  The \emph{leading term} of $\Phi(\Ob,\sigma)$ is
the pair
\[
(x_G,\lambda_{\rm lead}),\qquad
\lambda_{\rm lead}=\lambda_L-2\rho(\fu\cap\fs)+\rho(\fu),
\]
where $x_G$ is the image of $x_L$ in the KGB set of $G$.  Put
$\gamma_{\rm lead}(\sigma)=\dom((1+\theta_{x_G})\lambda_{\rm lead})$.
\end{definition}

This is the term with $A=\emptyset$ in the formula for $\Phi$ of \cite[Section~6]{KN}, coming
from the leading term of the $K$-type formula of $\sigma$, before
standardization.  It costs one call of
\at{twist\_by\_minus\_2rho\_u\_cap\_s}.  When the pair
$(x_G,\lambda_{\rm lead})$ is a standard parameter it is typically the lowest
term of $\Phi(\Ob,\sigma)$ (compare Proposition~\ref{p:cone}(3)).

The algorithm uses the leading term in three ways.
\begin{enumerate}
\item $\Phi(\Ob,\sigma)$ is computed only if $\gamma_{\rm lead}(\sigma)$ lies
in the region $R$.
\item In the first stages only those $\sigma$ are used for which
$(x_G,\lambda_{\rm lead})$ is standard.  For a principal orbit of an
equal rank group, where $L$ is a compact torus, this selects one Weyl chamber
out of $|W|$.  For a discrete series it leaves one or two $\sigma$.
\item It gives the range of $\sigma$ to enumerate.  Since
$\gamma\mapsto\langle\dom(\gamma),\ch\rho\rangle$ is convex and homogeneous,
\[
\|\sigma\|_{\ch\rho(G)}\le
\langle\gamma_{\rm lead}(\sigma),\ch\rho\rangle+c_\Ob,\qquad
c_\Ob=\max_{x_L}\big\langle\dom\big((1+\theta_{x_G})
(\rho(\fu)-2\rho(\fu\cap\fs))\big),\ch\rho\big\rangle .
\]
So if $N_R$ is the largest norm of a $K$-type in $R$, only the $\sigma$ with
$\|\sigma\|_{\ch\rho(G)}\le N_R+c_\Ob$ can pass test (1).  This replaces the
bound $2N$ of \at{Phi\_upto}; compare the conjecture in \cite[Section~8]{KN},
which asks for a bound of the form $cN$.  The natural bound is additive.
In Example~\ref{e:wrong} one has $N=3$ and $c_\Ob=6$, and the system becomes
solvable exactly when the $\sigma$ of norm $9=N+c_\Ob$ are included; the
bound $2N=6$ is too small.
\end{enumerate}

\section{The algorithm}
\label{s:alg}

\textbf{Input:} a real group $G$ and a list of parameters $p$ (for example a
block).  \textbf{Output:} for each $p$ the real orbits in $\AV(J(p))$ with
their multiplicities in the associated cycle.

\begin{enumerate}
\item Compute the real nilpotent orbits.  For each orbit $\Ob$ compute, when
it is first needed, the parabolic, the constant $c_\Ob$, the list of
$\sigma$ with $\|\sigma\|\le N_{\max}+c_\Ob$, and for each of them the
leading term.  The $\Phi(\Ob,\sigma)$ are computed on demand and stored, so
that each is computed at most once for the whole list of parameters.
\item For each $p$: let $w$ be \at{character\_formula(p)*0},
$d=2\,$\at{GK\_dim(p)}, and let the \emph{top} orbits be the real orbits of
dimension $d$.  Compute $\Gamma_{\max}(w)$, $\Gamma_{\min}(w)$.
\item Try the following column sets in turn, until $Tv=w$ has a rational
solution $v$.  In every stage a column $\Phi(\Ob,\sigma)$ is used only if all
its terms lie in the stated region.

\smallskip
\begin{tabular}{rlll}
stage & orbits & region & condition on the leading term\\
\midrule
0 & top & interval $I(w)$ & standard, and $\gamma_{\rm lead}\in\Gamma_{\min}(w)$\\
1 & top & interval $I(w)$ & standard, and in the region\\
2 & top & interval $I(w)$ & in the region\\
3 & top & cone $C(w)$ & in the region\\
4 & top & cone $C(w)$ & none\\
$5+j$ & dimension $\ge d_{j+1}$ & cone $C(w)$ & in the region
\end{tabular}

\smallskip\noindent
Here $d=d_0>d_1>d_2>\dots$ are the dimensions of the real orbits.  If no stage
succeeds, repeat stages $3,4,5,\dots$ with $C(w)$ replaced by $C_1(w)$, then
$C_2(w)$.
\item From the solution $v$ compute
$m_\Ob=\sum_\sigma v_{\Ob,\sigma}\dim\sigma$ for each top orbit.  Check that
these are nonnegative integers, not all zero.  The associated cycle is
$\sum m_\Ob\overline\Ob$.
\end{enumerate}

By Theorem~\ref{t:mult} the output is correct whenever step (3) succeeds and
the check in step (4) passes.  The rows of $T$ are the $K$-types occurring in
$w$ and in the columns used; nothing else needs to be enumerated.

\section{Experiments}
\label{s:exp}

The prototype of Appendix~\ref{s:code} was run on the block of the trivial
representation for the groups of Table~\ref{tab:old}.  Table~\ref{tab:blocks}
gives the number of $\Phi(\Ob,\sigma)$ computed, the size of the largest
linear system solved, the number of irreducibles solved at each stage, and
the time for the whole block.  For each group the old and the new time were
measured one after the other on the same machine.  (For $SU(3,2)$, where the
current code was stopped after 50 minutes, the prototype was stopped after
10 minutes without having finished.  Judging from Table~\ref{tab:single},
the time goes into listing the candidate $\sigma$: the current code lists
221514 of them for this group (Table~\ref{tab:old}), and the prototype lists
them in the same way; see Section~\ref{s:further}(1).  Presumably for the
same reason, a single large discrete series of $SU(3,3)$ did not finish in
10 minutes.)

\begin{table}[ht]
\centering
\begin{tabular}{lrrrrcrrrrr}
\toprule
 & & & \multicolumn{2}{c}{$\Phi$ computed} & largest
 & \multicolumn{3}{c}{solved at stage} & \multicolumn{2}{c}{time (s)}\\
$G$ & block & orbits & old & new & system & 0--2 & 3--4 & $\ge 5$ & old & new\\
\midrule
$SL(2,\R)$ & 3 & 3 & 17 & 5 & $2\times 2$ & 3 & 0 & 0 & 0.01 & 0.02\\
$SU(2,1)$ & 6 & 4 & 226 & 14 & $4\times 3$ & 6 & 0 & 0 & 0.2 & 0.1\\
$SL(3,\R)$ & 6 & 3 & 24 & 5 & $3\times 1$ & 6 & 0 & 0 & 0.03 & 0.02\\
$Sp(4,\R)$ & 12 & 8 & 522 & 53 & $25\times 34$ & 12 & 0 & 0 & 1.0 & 0.3\\
$G_2$ split & 12 & 6 & 558 & 55 & $16\times 40$ & 12 & 0 & 0 & 1.8 & 0.4\\
$SL(4,\R)$ & 15 & 7 & 824 & 78 & $21\times 24$ & 15 & 0 & 0 & 3.2 & 0.7\\
$SU(2,2)$ & 21 & 10 & 7395 & 357 & $30\times 10$ & 21 & 0 & 0 & 81 & 3.2\\
$SO(4,3)$ & 53 & 10 & 6564 & 980 & $52\times 212$ & 45 & 2 & 6 & 136 & 34\\
$SL(5,\R)$ & 55 & 7 & 1390 & 254 & $19\times 70$ & 41 & 2 & 12 & 68 & 11\\
$Sp(6,\R)$ & 53 & 19 & 29083 & 2035 & $132\times 188$ & 51 & 2 & 0 & $>3000$ & 51\\
\bottomrule
\end{tabular}
\caption{The new algorithm on the block of the trivial representation.
``Orbits'' is the number of real nilpotent orbits; ``largest system'' is the
size (rows $\times$ columns) of the largest linear system solved for an
irreducible in the block.  Stages 0--2 use the top orbits and the interval,
stages 3--4 the top orbits and the cone, and stages $\ge 5$ smaller orbits as
well (Section~\ref{s:alg}).  The times are for the whole block.}
\label{tab:blocks}
\end{table}

\mi\textbf{Agreement.}  The associated varieties were compared with the output
of the current code (\at{av}) for every irreducible of infinitesimal
character $\rho$ of $SL(2,\R)$, $SU(2,1)$, $SL(3,\R)$, $Sp(4,\R)$, $SO(4,1)$,
$Sp(1,1)$, $SU(3,1)$, split $G_2$, $SL(4,\R)$, $SL(2,\C)$, $SL(3,\C)$,
$Sp(4,\C)$, and for the block of the trivial representation of $SU(2,2)$,
$SO(4,3)$ and $SL(5,\R)$: 238 irreducibles in all.  They agree in
every case.  In every case the multiplicities are nonnegative integers.  They
were also compared at some singular and nonintegral infinitesimal characters
for $SL(2,\R)$, $SL(3,\R)$, $SU(2,1)$, $Sp(4,\R)$ and $G_2$; there the only
differences are those of Example~\ref{e:wrong}, where the current code is
wrong.

\mi\textbf{Multiplicities.}  The current code does not compute the associated
cycle.  Some values, at infinitesimal character $\rho$: for split $G_2$ the
irreducibles $J(x,\rho,\nu)$ with $x=3,7$ have multiplicity $3$, and those
with $x=5,6$ have multiplicity $2$, on a real form of the subregular orbit;
for $SL(4,\R)$, $J(x=12,\lambda=[3,3,2],\nu=\rho)$ has associated cycle
$3\,\overline{\Ob_2}+3\,\overline{\Ob_3}$; in the block of the trivial
representation the largest multiplicity is $6$ for $SO(4,3)$, $5$ for
$SL(5,\R)$ and $10$ for $Sp(6,\R)$.

\mi\textbf{Single representations.}  Table~\ref{tab:single} is for a single
irreducible of a larger group.  For a discrete series the linear system has
one or two unknowns, and one or two $\Phi(\Ob,\sigma)$ are computed.  Almost
all of the time is now spent in listing the candidate $\sigma$ and their
leading terms, which the prototype does in the same way as the current code
(all $\sigma$ up to a norm bound); see Section~\ref{s:further}(1).

\begin{table}[ht]
\centering
\begin{tabular}{lrrrcr}
\toprule
representation & orbits & $\sigma$ listed & $\Phi$ computed & system & time (s)\\
\midrule
$SO(4,3)$, large discrete series & 10 & 281 & 1 & $1\times 1$ & 0.2\\
$Sp(6,\R)$, large discrete series & 19 & 1250 & 2 & $2\times 2$ & 0.5\\
$Sp(6,\R)$, trivial & 19 & 146 & 7 & $60\times 1$ & 0.3\\
$SO(5,4)$, large discrete series & 21 & 5593 & 1 & $1\times 1$ & 7.2\\
$Sp(8,\R)$, large discrete series & 41 & 24498 & 2 & $2\times 2$ & 56\\
$Sp(8,\R)$, trivial & 41 & 973 & 9 & $285\times 1$ & 5.3\\
$F_4$ split, trivial & 27 & 606 & 3 & $290\times 1$ & 5.1\\
$F_4$ split, large discrete series & 27 & 28369 & 1 & $1\times 1$ & 106\\
\bottomrule
\end{tabular}
\caption{The new algorithm on single representations.  All are solved at
stage 0.}
\label{tab:single}
\end{table}

\section{Further improvements}
\label{s:further}

\begin{enumerate}
\item \emph{Enumerate $\sigma$ from the region.}  The prototype lists all
$\sigma$ with $\|\sigma\|\le N_R+c_\Ob$ and then tests the leading term.  It
would be much better to go the other way: list the dominant weights $\gamma$
in the region $R$, for each one the $L$-dominant $W$-conjugates $\gamma'$, and
solve $(1+\theta_{x_L})(\lambda_L-2\rho(\fu\cap\fs)+\rho(\fu))=\gamma'$ for
$\lambda_L$.  This is what \at{lambda\_rho\_from\_dominant\_height\_upto} does,
with the list of $\gamma$ given by the region instead of a norm bound.  For a
discrete series of $F_4$ this lists at most $|W|=1152$ candidates instead of
28369.
\item \emph{Speed up $\Phi$.}  In \at{Phi}, the list \at{rho\_shifts(O,x)}
depends only on the orbit and on $x$, but it is recomputed for every term of
every $\sigma$ (including a call of \at{dominant(O)}).  The prototype
computes it once for each $x$.  The same applies to
\at{parabolic\_of\_orbit(O)} and its Levi factor.  For the zero orbit,
$\Phi(0,\sigma)$ is just \at{K\_type\_formula}.
\item \emph{Closure order.}  When the top orbits do not suffice
(Example~\ref{e:so43}) the prototype adds all orbits of the next dimension.
With the closure order on real orbits one would add only the maximal orbits
in the intersection of the closures of two top orbits.
\item \emph{Reductive groups.}  The region $R$ is finite without a choice of
$v$.  With the enumeration of (1) the vector $v$ of \cite[Section~8]{KN} would
disappear from the algorithm.
\item \emph{Translation.}  Translation between two regular infinitesimal
characters in the same chamber is an equivalence and does not change the
associated variety.  At infinitesimal character $\rho$ the cone never had to
be enlarged (Example~\ref{e:sl2} shows that elsewhere it must be), and the
norms are smallest, so for a regular parameter far from the walls it may be
cheaper to translate first.  This gives the associated variety but not the
multiplicities.
\item \emph{What should be proved.}  Observation~\ref{o:branch}; a precise
form of Working hypothesis~\ref{h:interval} with the cone $C_k(w)$ and an explicit
$k$; and a description of the image of $\mu_*$ that explains
Example~\ref{e:so43}.
\end{enumerate}

\appendix
\section{The prototype}
\label{s:code}

The following is the file \at{K\_Nilpotent\_improved.at} with which the
experiments were done.  It uses \at{K\_Nilpotent.at} for $\Phi$ and for the
enumeration of $L\cap K$-types.  The main functions are
\at{kf\_associated\_cycles(G,params)} and \at{kf\_associated\_cycle(p)}.  For
example:
\begin{verbatim}
atlas> set G=Sp(4,R)
atlas> kf_associated_cycle(parameter(KGB(G,9),[2,1],[3,3]/2))
Value: [(([ 1, 1 ],KGB element #0,()),2)]
\end{verbatim}

{\footnotesize
\begin{verbatim}
<K_Nilpotent.at
<associated_variety_annihilator.at { for GK_dim@Param }

{ Prototype of the improved K-nilpotent algorithm; see knilpotentImproved.tex.
  Main functions:
    kf_associated_cycles(G,[Param])  associated cycles of a list of irreducibles
    kf_associated_cycle(Param)       associated cycle of one irreducible
}

set kf_verbose=false
set kf_max_margin=2

{ gamma(mu): the dominant weight (1+theta_x)lambda attached to a K-type }
set kf_gamma (KType t) = ratvec:
  dominant(t.root_datum,(1+involution(t.x))*t.lambda)

{ dominance order: g <= L iff L-g is a nonnegative combination of simple roots }
set kf_leq (RootDatum rd) = (ratvec,ratvec->bool):
  let fcw=fundamental_coweights(rd), rb=radical_basis(rd) in
  (ratvec g,ratvec L) bool:
    let diff=L-g in all(for cw in fcw do diff*cw>=0 od) and =(diff*rb)

set kf_maximal (RootDatum rd,[ratvec] S) = [ratvec]:
  let leq=kf_leq(rd) in
  ##for g@i in S do
     if any(for h@j in S do j!=i and leq(g,h) and (not leq(h,g) or j<i) od)
     then [] else [g] fi od

set kf_minimal (RootDatum rd,[ratvec] S) = [ratvec]:
  let leq=kf_leq(rd) in
  ##for g@i in S do
     if any(for h@j in S do j!=i and leq(h,g) and (not leq(g,h) or j<i) od)
     then [] else [g] fi od

{ the leading raw term of Phi(O,sigma): the parameter obtained from the STFL
  parameter of sigma by the twist and theta-induction, before standardization.
  Returns its gamma, and whether it is already a standard parameter of G }
set kf_lead (RealForm G,Parabolic P,ratvec shift_u,KType mu_L) = (ratvec,bool):
  let q=twist_by_minus_2rho_u_cap_s(P,mu_L.param)
  then x_G=embed_KGB(x(q),G), lam=lambda(q)+shift_u
  in ( dominant(G.root_datum,(1+involution(x_G))*lam)
     , is_standard(parameter(x_G,lam,null(rank(G)))) )

{ the constant c_O: if the leading term of Phi(O,sigma) has norm <= N then
  sigma has G-norm <= N + c_O }
set kf_shift_bound (RealForm G,Parabolic P) = int:
  let L=Levi(P), rd=G.root_datum, v=rho_check(G) then shift=rho(G)-rho(L) in
  max(0)(for x_L in KGB(L) do
     let x_G=embed_KGB(x_L,G)
     then s=(1+involution(x_G))*(shift-two_rho_u_cap_s(P,x_G))
     in ceil(dominant(rd,s)*v) od)

{ the list of (sum of A, #A), A a subset of the restricted roots of s[1] }
set kf_rho_shifts (RealNilpotent O_dom,KGBElt x_G) = [(vec,int)]:
  let r=rank(O_dom.root_datum) in
  [(vec:null(r),0)] ##
  for A in subsets_of_s_one_nc_cx_restricted_roots(O_dom,x_G)
  do (sum(r # A),#A) od

{ same as Phi@(RealNilpotent,KType) of K_Nilpotent.at, with the data that
  depends only on the orbit precomputed }
set kf_Phi (RealForm G,Parabolic P,[[(vec,int)]] shifts,KType mu_L) = KTypePol:
  let Ktw=twist_by_minus_2rho_u_cap_s(P,K_type_formula(mu_L))
  , result=null_K_module(G) in
  for c@p_L in Ktw do
    for (shift,j) in shifts[#x(p_L)] do
      let q=K_type(x(p_L),lambda_rho(p_L)+shift) in
      result+:=minus_1^j*c*theta_induce_standard(param(q),G).K_type_pol
    od
  od; result

{ solve sum a(i,sigma) Phi(O_i,sigma) = P over Q; cols[i] are the columns of
  the i-th orbit. Returns m_i = sum_sigma a(i,sigma) dim(sigma), i<ntop }
set kf_solve ([[(KType,KTypePol)]] cols,int ntop,KTypePol P) =
      (bool,[rat],int,int):
  let flat=##for c in cols do for (,Q) in c do Q od od in
  if #flat=0 then (=P,for i:ntop do rat:0 od,#P,0) else
  let basis=monomials(flat#P)
  then M=mat: for Q in flat do Vector(basis,Q) od
  then v0=solve(M,ratvec: Vector(basis,P)) in
  if not succeeds(v0) then (false,[rat]:[],#basis,#flat) else
  let v=requisition(v0), offs=0 in
  ( true
  , for c@i in cols do
      let m=rat:0 in
      for (s,)@j in c do m+:=v[offs+j]*dimension(s) od; offs+:=#c;
      if i<ntop then [m] else [] fi
    od.##
  , #basis, #flat) fi fi

{ result for one parameter: solved?, indices of the orbits of dimension
  2*GK_dim, their multiplicities, stage at which a solution was found,
  size of the linear system }
set_type KF_result =
  (bool solved,[int] top,[rat] mults,int stage,int ncols,int nrows)

{ one pass, with the cone enlarged by margin*2rho }
set kf_associated_cycles_margin \
   ([RealNilpotent] orbits,RealForm G,[Param] params,int margin) =
      ([KF_result],[int]):
  let rd=G.root_datum, v=rho_check(G)
  then extra=margin*two_rho(rd)
  then n=#orbits, leq=kf_leq(rd)
  then Ps=for p in params do character_formula(p).K_type_pol od
  then Nmax=max(0)(for P in Ps do K_norm(P) od)+rat_as_int(extra*v)
  { per-orbit caches }
  then ready=for i:n do false od
  , pars=for O in orbits do parabolic_of_orbit(O) od
  , sigs=for i:n do [KType]:[] od
  , leads=for i:n do [ratvec]:[] od
  , leadstd=for i:n do [bool]:[] od
  , shifts=for i:n do [[(vec,int)]]:[] od
  , phis=for i:n do [KTypePol]:[] od
  , phig=for i:n do [[ratvec]]:[] od  { gammas of the terms of each Phi }
  , done=for i:n do [bool]:[] od
  , nphi=0, nenum=0
  then prepare(int i)=void:
    if not ready[i] then
      let P=pars[i] then L=Levi(P), O_dom=dominant(orbits[i])
      then shift_u=rho(G)-rho(L)
      then ss=K_parameters_norm_upto(L,G,Nmax+kf_shift_bound(G,P),v)
      then ld=for s in ss do kf_lead(G,P,shift_u,s) od in
      sigs[i]:=ss; nenum+:=#ss;
      leads[i]:=for (g,) in ld do g od;
      leadstd[i]:=for (,b) in ld do b od;
      shifts[i]:=for x_L in KGB(L) do kf_rho_shifts(O_dom,embed_KGB(x_L,G)) od;
      phis[i]:=for s in ss do null_K_module(G) od;
      phig[i]:=for s in ss do [ratvec]:[] od;
      done[i]:=for s in ss do false od;
      ready[i]:=true
    fi
  { columns of orbit i: those sigma whose leading term satisfies lead_ok (and
    is standard, if std_only) and all of whose terms satisfy term_ok }
  then columns(int i,bool std_only,(ratvec->bool) lead_ok,
               (ratvec->bool) term_ok) = [(KType,KTypePol)]:
    prepare(i);
    let ph=phis[i], pg=phig[i], dn=done[i], ld=leads[i], lstd=leadstd[i]
    then cols=
    ##for s@k in sigs[i] do
       if (std_only and not lstd[k]) or not lead_ok(ld[k]) then [] else
         if not dn[k] then
           let Q=kf_Phi(G,pars[i],shifts[i],s) in
           ph[k]:=Q; pg[k]:=for @t in Q do kf_gamma(t) od; dn[k]:=true;
           nphi+:=1
         fi;
         if #pg[k]>0 and all(for g in pg[k] do term_ok(g) od)
         then [(s,ph[k])] else [] fi
       fi od
    in phis[i]:=ph; phig[i]:=pg; done[i]:=dn; cols
  then results=
  for p@ip in params do
    let P=Ps[ip] then dm=2*GK_dim(p)
    then top=##for O@i in orbits do if dimension(O)=dm then [i] else [] fi od
    then lowdims=
      sort_u(##for O in orbits
               do if dimension(O)<dm then [dimension(O)] else [] fi od).reverse
    then gs=for @t in P do kf_gamma(t) od
    then tops=kf_maximal(rd,gs), bots=kf_minimal(rd,gs)
    then below(ratvec g)=bool: any(for L in tops do leq(g,L+extra) od)
    then between(ratvec g)=bool: below(g) and any(for L in bots do leq(L,g) od)
    , lowest(ratvec g)=bool: any(for L in bots do g=L od)
    , always(ratvec g)=bool: true
    then res=KF_result:(false,top,[rat]:[],-1,0,0)
    then try(int stage,[int] orbs,bool std,(ratvec->bool) lead_ok,
             (ratvec->bool) term_ok)=void:
      if not res.solved then
        let (ok,m,nr,nc)=
          kf_solve(for i in orbs do columns(i,std,lead_ok,term_ok) od,#top,P)
        in if ok then res:=(true,top,m,stage+20*margin,nc,nr) fi
      fi
    in
    if margin=0 then
      try(0,top,true,lowest,between);   { leading term a lowest term of P }
      try(1,top,true,between,between);  { interval, standard leading term }
      try(2,top,false,between,between)  { interval }
    fi;
    try(3,top,false,below,below);       { cone }
    try(4,top,false,always,below);      { cone, any leading term }
    let use=top in                      { add smaller orbits }
    for dl@k in lowdims do
      use##:=##for O@i in orbits do if dimension(O)=dl then [i] else [] fi od;
      try(5+k,use,false,below,below)
    od;
    if kf_verbose then
      prints("kf: parameter ",ip," stage ",res.stage," system ",
             res.nrows,"x",res.ncols) fi;
    res
  od
  in (results,[nenum,nphi])

{ driver: margin 0 first; parameters that remain unsolved are retried with
  larger margins. Returns the real orbits, one KF_result per parameter, and
  [number of sigma enumerated, number of Phi(O,sigma) computed] }
set kf_associated_cycles (RealForm G,[Param] params) =
      ([RealNilpotent],[KF_result],[int]):
  let orbits=real_nilpotent_orbits(G)
  then (results,stats)=kf_associated_cycles_margin(orbits,G,params,0)
  , margin=1 in
  while margin<=kf_max_margin and any(for r in results do not r.solved od) do
    let todo=##for r@k in results do if r.solved then [] else [k] fi od
    then (res2,st2)=
      kf_associated_cycles_margin(orbits,G,for k in todo do params[k] od,margin)
    in for k@j in todo do results[k]:=res2[j] od;
    stats:=[stats[0]+st2[0],stats[1]+st2[1]]; margin+:=1
  od; (orbits,results,stats)

{ associated cycle of one irreducible: the orbits in AV(J(p)), each with its
  multiplicity }
set kf_associated_cycle (Param p) = [(RealNilpotent,int)]:
  let (orbits,res,)=kf_associated_cycles(p.real_form,[p]) then r=res[0] in
  assert(r.solved,"no solution found; increase kf_max_margin");
  assert(any(for m in r.mults do !=m od),"all multiplicities vanish");
  ##for i@k in r.top do
     if !=r.mults[k] then [(orbits[i],rat_as_int(r.mults[k]))] else [] fi od
\end{verbatim}
}

\begin{thebibliography}{9}

\bibitem[AV]{AV}
J.~Adams and D.~Vogan, \emph{Associated varieties for real reductive
groups}, arXiv:2103.11836.

\bibitem[KN]{KN}
J.~Adams, \emph{$K$-nilpotent algorithm}, notes (\at{knilpotent.tex}).

\bibitem[V]{Vogan91}
D.~Vogan, \emph{Associated varieties and unipotent representations}, in:
Harmonic analysis on reductive groups (Brunswick, ME, 1989), Progr.\ Math.\
\textbf{101}, Birkh\"auser, 1991, 315--388.

\end{thebibliography}

\end{document}
